% 05_efficiency_tips.tex — Efficiency Tips for Engineers: Insights from Mo Tiwani

\section{Efficiency Tips for Engineers: Insights from Mo Tiwani}
\label{sec:efficiency_tips}

\andrea{wonder if we really need separate ``key takeaways'' sections, since even the prose sections are quite concise --- about the same number of sentences in each}

\subsection{Background and Perspectives}

Mo Tiwani's journey to obtaining his PhD was unconventional, as he spent multiple years in industry before returning to academia.
Reflecting on his career, he emphasizes the importance of intentional time management.
Engineers must identify their goals and ensure their activities align with these objectives.
Empirical tools such as RescueTime can be invaluable for tracking and optimizing time usage, and highlight how time is spent based on data, not hindsight.
Optimal productivity levels are typically around 80\%, with 65\% being a satisfactory benchmark.

\subsection{Strategic Work Tips}

Efficiency in engineering often has a positive correlation with experience.
Senior engineers tend to use open frameworks and avoid getting lost in details of full re-implementation, operating more efficiently by ethically reusing work rather than redoing it.
A research related example comes in the form of literature reviews.
For literature reviews, engineers can employ the ``Nearest Neighbor'' method---leveraging existing reviews while adding new insights.
In mathematical contexts, understanding theorem statements can often suffice, rather than delving deeply into the proofs.
Of course, in an educational setting, or if the task is something which you need to understand in the future at a deeper level, then full implementation is the best path.
Knowing the difference between when to focus deeply for recurrent knowledge, and optimizing via existing work for a one-off is a key distinction.
Building on this, pre-deciding important matters streamlines decision-making in general.
One example is the Eisenhower Matrix.
The Eisenhower Matrix \textcolor{red}{(filler example below, perhaps we make a custom one for the final draft)} is a useful tool for distinguishing between urgent and important tasks, facilitating proper time allocation.
\begin{figure}[htbp]
  \centering
  \includegraphics[width=0.4\textwidth]{Eisenhower_Matrix_WikimediaCommons_Filler.png}
  \caption{Eisenhower Matrix. Cmglee, via Wikimedia Commons}
  \label{fig:eisenhower}
\end{figure}
Importantly, striving for perfection can hinder progress, and it is essential to balance high standards with timely execution.

\subsection{Engineering vs.\ Research}

There is a crucial distinction between engineering and research workflows.
In engineering, both the problem and solution are usually well-defined, facilitating efficiency through established methodologies.
Conversely, research environments often feature unknown problems and solutions, necessitating rapid iteration and tolerance for failure.
For a researcher, the largest cost is the cost of failure.
On the other hand, engineers should maximize the value-to-time ratio, whereas researchers should focus on failing fast to gain insights.
The central motivating example covered by Tiwani, which he originally heard from Google X, is the monkey on a pedestal.
Imagine you are given the following task,
\begin{quote}
\small
``You want to train a monkey to recite Shakespeare on top of a pedestal.''
\end{quote}
An engineer may approach this problem by first building a pedestal, since it is something that they know how to do.
By contrast, a researcher may start with the highest uncertainty task, which is to train the monkey to recite Shakespeare.
Of course, both paths will eventually lead to failure, since it is impossible to teach a monkey to recite Shakespeare, but the researcher will have saved time that the engineer built the pedestal with, since they identified the highest uncertainty task first.
Both frameworks are valuable, and knowing when to use each is perhaps most important.

\subsection{Motivation and Procrastination}

Motivation, according to Tiwani, \textit{follows} action rather than preceding it.
Initiating tasks, even minimally, generates momentum.
In fact, tackling the hardest tasks first can simplify the rest of the workday.
As Tiwani puts it, jumping into cold water is difficult, but waiting doesn't make it any warmer!
It is important to note that procrastination is not a function of laziness but rather of perceived effort and emotional management, there are real tradeoffs to dedicating time to a task, but sometimes we tend to overestimate these.
Lowering the activation energy required to start tasks and boosting accountability through peer groups represent effective strategies to combat procrastination.
Lastly, passion for the work you are doing significantly influences engagement and productivity.

\subsection{Operations and Focus}

Integrating breaks is essential for sustained productivity.
Techniques like dopamine fasting---allowing periods of boredom such as lying on one's bed for 15 min---can enhance focus and creativity.
Another strategy to increase productivity is time blocking.
Time blocking prevents decision paralysis and ensures generous allocation to avoid stress from falling behind.
It is better to err on the side of allocating too much time for a task, and build momentum by completing it with ample time than to underestimate what is required and consistently fall behind.
Strengths should be leveraged for energy, while weaknesses should be managed carefully.
The Pomodoro Technique ensures breaks are genuine and restorative.
Focus, particularly deep work, can triple efficiency.
Managing distractions from smartphones through techniques like grayscale mode and occasional physical separation, can sustain focus.

\subsection{Tactical Implementation and Efficiency Enhancements}

Frequent tasks should be optimized for efficiency using hotkeys and multiple screens where feasible.
Identify and remove imaginary constraints that may unnecessarily dictate task sequences.
Memorize important information when possible and use printed materials for enhanced focus.
Tasks that take less than two minutes should be addressed immediately.
Strategic declination of tasks that do not align with primary objectives is essential.
\andrea{more on the prose/takeaways distinction (or lack thereof) --- e.g., these sections are almost identical}

\subsection{Managing Issues and Learning}

When facing issues such as missed deadlines, immediate communication with affected parties and providing updated deadlines is crucial.
Mistakes are learning opportunities, and systematic reuse of work can enhance efficiency.
Seek advice from experienced individuals to mitigate risks and appreciate the process rather than focusing solely on the outcome.
